\section{Related work}
\citet{bouthillier2021} derive variance inflation from correlated trainings on one task, whereas we account for covariance across benchmarks within a replicate run. \citet{reimers2017}, \citet{zhou2020}, \citet{madaan2024}, and \citet{fehlauer2025} study seed-related variability, and \citet{dodge2020} test whether good initialisations transfer across tasks. \citet{jordan2024} finds run deviations in accuracy nearly independent across ImageNet evaluation sets, which our accuracy interval allows, while our margin covariance sits within formats. \citet{miller2024} gives clustered and paired item-sampling errors with run variance fixed, and \citet{henderson2018} and \citet{agarwal2021} resample runs within a task, which discards the paired run effects we keep. We know of no study that estimates the covariance of run noise between benchmarks or an effective benchmark count from that covariance.

Our cross-half estimator adapts the split-half construction of \citet{spearman1910,brown1910} and the variance-component ideas of \citet{spearman1904}, \citet{cronbach1951}, and \citet{searle1992}, and what we add is the estimand and its measurement. \citet{cheverud1988} motivates our phenotypic plug-in comparison (Table~\ref{tab:prediction}), and we borrow the covariance analogy of \citet{martin1977} without its variance components.

Our few-cluster inference follows \citet{cameron2008} and \citet{mackinnon2017}. \citet{heineman2025} relate per-benchmark signal and noise on DataDecide and OLMo checkpoints to decision accuracy without estimating covariance between benchmarks. Across models, benchmark scores share a low-dimensional capability structure \citep{ruan2024,burnell2023revealing}, whereas we measure covariance of run deviations with the configuration held fixed.

\section{Discussion and limitations}
We read the accuracy result as a bound rather than as evidence of independence. The full-battery accuracy interval includes one at a design that would detect a true ratio of 1.10 in 0.135 of replicates, and the primary analysis gives an upper limit of 1.157 at the scored checkpoint. Exploratory cuts whose intervals exclude one, without BoolQ at 1.558, at the adjacent checkpoint and on the held-out tasks, don't settle it.

Apart from the registered test, which failed at its registered scope, and the two PolyPythias rules fixed before scoring, our study remains exploratory. We didn't score the 375 DataDecide runs on the second bank, which leaves the 1.244 untested on fresh items and the shared-item check unrun in the family that supplies our headline. Our populations stop at 1B and at likelihood-scored multiple choice, since no larger public multi-seed population shares replicate item scores (Appendix~\ref{app:design}) and generative tasks lack our margin. The passage-aware split moves both estimates within re-split variation, a null that reflects BoolQ's passage structure more than a sensitive test. A scalar item effect shared across benchmarks and repeated across halves would enter the numerator whole, and shares of 0.124 and 0.009 of item-noise variance reproduce 1.244 and 1.078, an item effect we haven't measured or separated from seed covariance (Appendix~\ref{app:calibration}).\outcome{\Rtwocase}{1}{R2 not supported}{ The cross-bank check in PolyPythias constrains that item effect in one family, and it can't reach an item effect that the two banks share.} A scalar component shared by every benchmark would also enter the cross-format pairs, whose diagnostic of 0.921 constrains a uniform component but not one that follows scoring format. At a share of 0.05 only 0.007 of simulated replicates reach the observed margin value.

We cluster on recipes, but the release shares one seed label across every recipe inside a size band, and a run effect following the band, which seed-matched paired comparisons cancel, would break that choice and take simulated margin coverage to 0.821 at a quarter of latent variance and 0.580 at a half. A permutation test that aligns runs by batch slot bounds that effect at 0.044 of the seed covariance on margins, where simulated coverage is 0.941 at a share of 0.05, and that test's power is only 0.490 at the bound. At 1B alone the cluster-$t$ margin interval of 0.811 to 1.628 includes one (Appendix~\ref{app:supplementary}).

We recommend paired replicate scores on the chosen battery, which give the aggregate seed standard deviation directly, and a benchmark-removal check. Out of sample on margins, independence predicts a configuration's measured ratio with a mean absolute log error of 0.339 against 0.536 for the plug-in and 0.512 for our decomposition (Appendix~\ref{app:transport}). Our ratio therefore measures this battery and helps only where replicates are missing and the target battery's own ratio is close. Without replicates, the format fit of Appendix~\ref{app:formatfit} gives a fallback scaling we haven't tested outside this battery, and seed-matched paired comparisons need little correction, with median paired-difference ratio 1.024 (Appendix~\ref{app:paired}). Scaling the independence standard error by a ratio from the same margin population returns the simulated false-positive rate from 0.113 to 0.051, and a separate transfer simulation to a battery whose own ratio is 1.5 leaves it at 0.101. Under equicorrelated seed covariance at 125 configurations, three runs each let the margin interval exclude one in 0.997 of replicates, and accuracy needs eight to reach 0.803.

\label{maintext:end}
\section*{AI use statement}
We used generative AI tools for three purposes. They helped polish the wording of the manuscript, they helped us find related work, and they gave minor help with analysis code. The authors designed the study, wrote the paper, checked every AI suggestion before using it, and take full responsibility for the code, the citations and the final text.

\section*{Ethics statement}
This study analyses released model outputs and introduces no new participant recruitment or data collection. Its intended use is to improve uncertainty reporting in model evaluation. An independence approximation used outside the studied battery could understate uncertainty, which keeps the recommendations conditional on the measured population and score definition. Use of derived artifacts must retain the applicable source attribution and licensing requirements.

\section*{Reproducibility statement}
Section~\ref{sec:estimator} and Appendix~\ref{app:derivations} define the estimator and its assumptions. Appendix~\ref{app:implementation} specifies the interval calculation and data reduction. Appendix~\ref{app:design} records the populations and planned checks. Appendix~\ref{app:supplementary} reports later sensitivity analyses alongside the primary results. The released archive also holds an extended supplement with the full-length versions of the compressed appendix sections. The primary calculation uses 4,999 bootstrap draws, 2,000 calibration replicates, 500 permutation replicates, 50 item re-splits, and master seed 20260101, and reproduction requires the per-item reductions, item ordering and half assignments, selected checkpoint identifiers, and the corresponding analysis scripts. We release all four of them alongside the paper. Neither the analysis plan nor the protocol and rule commits carry an independently corroborated timestamp, and the plan's screening split was never formed. A separate implementation run on Kaggle reproduces the primary estimates, the BoolQ-removal estimates, and their intervals to three decimals from the released per-item scores, a run that also produced the passage-aware split, the coverage simulation, the operational prediction comparison, and the seed-matched pair analysis. The PolyPythias transfer, five-size rescore, and checkpoint results come from earlier execution records that this revision didn't rerun, and the loss-proxy rescoring wasn't run. The registered test ships with its protocol, the dated amendments that record each change to it, the frozen hash of its request file, and the kernels that scored and analysed it, and the last twenty runs at 1B were scored off Kaggle on three RTX 3090 cards. The PolyPythias battery and the second bank were scored in float32 with transformers 4.57.1 on Kaggle T4 cards, each run checking its cached scores against an uncached reference on 24 items to within 0.001 nats, and the release ships the frozen hashes of both request files with the kernels that scored them. Separate CPU kernels on the released runs produced the recipe-decision, seed-count, battery-size, and item-count analyses.

